\chapter{Learning Tasks on Graphs}
\label{ch:applications}
\minitoc

Once a domain has been represented as graph-structured data, a natural
question arises: what machine learning tasks can be defined on a graph?
A graph can be viewed at three levels: its nodes, its edges, and the graph
as a whole. Each gives rise to a family of prediction tasks: predicting a
property of each node, predicting the existence or type of a connection
between pairs of nodes, or predicting a property of the entire graph. These are known,
respectively, as \textit{node-level}, \textit{edge-level}, and
\textit{graph-level tasks} (Figure~\ref{fig:graph-learning-tasks}), and they structure the 
architectures and implementations developed in this work.

\begin{figure}[!h]
  \centering
  \includegraphics[width=\textwidth]{graph-learning-tasks}
  \caption[The three levels of graph learning tasks.]{%
    The three standard levels at which predictions are commonly made on a graph. 
    A node-level task predicts a property of an individual node (highlighted); an
    edge-level task predicts a property of a connection between two nodes, such
    as whether it exists (dashed); and a graph-level task predicts a single
    property of the entire graph (framed). 
  }
  \label{fig:graph-learning-tasks}
\end{figure}

These three levels correspond to the standard categories of graph learning tasks, 
but they do not exhaust the field. Other settings include \textit{subgraph-level tasks}, 
which focus on groups of related nodes, and \textit{graph generation tasks}, 
which aim to generate new graphs, such as candidate molecular structures with 
desirable properties. Although these settings rely on the same representational 
principles introduced here, they lie outside the scope of this work, 
which focuses on the three standard task levels.

The following sections examine each of these levels
in turn and present representative applications that motivate them.

\section{Node-Level Tasks}
\label{sec:node-level}

A \textit{node-level task} asks for a prediction about each individual node of a
graph, drawing on both the node's own attributes and its position within the
surrounding structure. The most common instance is \textit{node
classification}: assigning each node to one of a set of categories. 
A representative setting is a citation network, in
which the nodes are scientific papers, the edges are citations between them, and
the goal is to infer the topic of each paper (Figure~\ref{fig:node-level}). A paper's own 
text is informative, but so is its neighborhood: papers tend to cite others on related subjects,
so the topics of a paper's neighbors are strong evidence for its own. Node
classification exploits exactly this, combining what a node contains with what
surrounds it. 

\begin{figure*}[t]
  \includegraphics[width=\textwidth]{node-level-tasks}
  \caption[Node-level tasks.]{%
    Left: the node classification task in the abstract. Some nodes carry a known
    label (shown by color) while others are unknown (marked with a question mark).
    Right: a concrete instance as a citation network,
    where each node is a scientific paper and the predicted label is its topic.
  }
  \label{fig:node-level}
\end{figure*}

The target need not be a category. A graph representation of a protein, for example, 
can be constructed where nodes represent amino acid residues and edges connect 
residues that are close in space. One may predict a discrete property of each
residue, such as its functional role, or a continuous property, such as its position
in three dimensions, a formulation that underlies recent advances in protein
structure prediction~\cite[-16pc]{jumper2021}. 

\newpage

The same framing extends to economic settings. In a financial network whose
nodes are accounts and whose edges are transactions between them, fraud
detection amounts to flagging each account as legitimate or fraudulent: a node
classification problem in which fraudulent behavior is often more evident in the
surrounding network of connections than in any single account considered in isolation.

\section{Edge-Level Tasks}
\label{sec:edge-level}

\begin{figure*}[t]
  \centering
  \includegraphics[width=\textwidth]{edge-level-tasks}
  \caption[Edge-level tasks.]{%
    Left: the link prediction task in the abstract. Solid edges are observed
    connections; dashed edges marked with a question mark are candidate
    connections whose existence the model must predict. Right: a concrete
    instance as a recommendation setting, where the two kinds of entity (users,
    drawn as circles, and products, drawn as squares) are joined by observed
    interactions, and the task is to predict further user--product links.
  }
  \label{fig:edge-level}
\end{figure*}

An \textit{edge-level task} focuses on relationships between entities. 
The central example is \textit{link prediction}: given the edges already present 
in a graph, predict which additional edges are likely to exist. Recommendation is one 
of its most familiar applications (Figure~\ref{fig:edge-level}). 
In a product network, where nodes represent products and edges indicate co-purchases,
 predicting missing edges identifies items likely to be associated even if 
they have not yet been observed together. A richer formulation introduces 
users and products as distinct node types, with edges representing interactions 
such as purchases or ratings; link prediction then becomes the recommendation 
problem itself, anticipating which products a user is likely to prefer next.

What makes the graph formulation effective is that it learns from the relational
structure rather than treating each entity in isolation. A recommender system that treats
products only by their individual attributes, ignoring the network of interactions
around them, discards precisely the signal that reveals which items behave
similarly; methods that propagate information along the graph structure can outperform 
approaches that rely only on individual attributes, a principle exploited by 
graph-based recommender systems deployed at web scale~\cite[-7.5pc]{ying2018}.

Link prediction extends well beyond commerce. Given a pair of drugs, one may
predict whether taking them together produces an adverse interaction, by placing
both within a larger biological network and asking whether a particular kind of
connection should join them~\cite[4pc]{zitnik2018}. The same formulation is also used to
complete \textit{knowledge graphs}, large networks of facts in which entities are connected
by relations: many true relations are simply unobserved, and link
prediction proposes the most plausible ones.

\section{Graph-Level Tasks}
\label{sec:graph-level}

\begin{figure*}[t]
  \centering
  \includegraphics[width=\textwidth]{graph-level-tasks}
  \caption[Graph-level tasks.]{%
    Left: the graph classification task in the abstract. The data is a
    collection of distinct graphs, differing in size and shape, each assigned a
    single label for the graph as a whole. Right: a concrete instance as
    molecular property prediction, where each graph is a molecule (atoms as
    nodes, bonds as edges) labeled by a property such as being active or
    inactive against a target.
  }
  \label{fig:graph-level}
\end{figure*}

A \textit{graph-level task} treats an entire graph as a single example and
predicts a property of the whole. The representative task is \textit{graph
classification}: assigning a label to a complete graph
(Figure~\ref{fig:graph-level}). Here the data is not one large network but a
collection of many separate graphs, each with its own label, and the model must
summarize the structure and attributes of a graph into a single
prediction. Molecular property prediction is the canonical example. A molecule
is naturally a graph, with atoms as nodes and chemical bonds as edges, and
prediction tasks about it are graph-level: whether a compound is toxic, or whether it is
active against a given biological target. Because molecules vary in size and
structure, this regime calls for a model able to handle graphs
that vary in this way and return one prediction per graph.

This capability feeds directly into drug discovery, where large libraries of
candidate compounds are screened to prioritize the few worth synthesizing and
testing in the laboratory. A graph-level model trained to predict antibacterial
activity, for instance, identified halicin, a compound structurally unlike
conventional antibiotics and effective against resistant
bacteria~\cite[-4.5pc]{stokes2020}. The same idea applies in materials science, where the
atomic structure of a material is represented as a graph and a graph-level model
predicts a physical property such as its stability or its mechanical strength.

Across all three levels, the graph is used in different
ways, with the level of the task determining what is predicted and how a model
must process the structure to predict it.